\documentclass[11pt,a4paper]{article}

\usepackage[utf8]{inputenc}
\usepackage[T1]{fontenc}
\usepackage[margin=2.5cm]{geometry}
\usepackage{amsmath,amssymb}
\usepackage{booktabs}
\usepackage{graphicx}
\usepackage{hyperref}
\usepackage{microtype}
\usepackage{parskip}
\usepackage{xcolor}
\usepackage{caption}
\usepackage{subcaption}

\hypersetup{
  colorlinks=true,
  linkcolor=blue!60!black,
  citecolor=blue!60!black,
  urlcolor=blue!60!black,
}

\title{\textbf{Volunteer-Based Last-Mile Delivery over Metro Networks}\\
       \large A Discrete-Event Simulation Study}
\author{João Pedro Pedroso\\
        \small Faculty of Sciences, University of Porto\\
        \small \texttt{jpp@fc.up.pt}}
\date{May 2026}

\begin{document}

\maketitle

\begin{abstract}
We study volunteer-based last-mile parcel delivery (VLMD), a system in which
ordinary commuters carry parcels between station lockers on their regular metro
trips. Using a discrete-event simulator calibrated to the Porto Metro network
(85 stations, $\sim$195{,}000 daily OD trips), we measure delivery rate and
delivery span as a function of (i) the fraction of passengers acting as
volunteers, (ii) the spatial correlation between parcel and passenger OD
patterns, and (iii) the transshipment heuristic. Four heuristics are evaluated:
\texttt{hop} and \texttt{flow} (monotone-progress Phase~1 only) and
\texttt{hop+} and \texttt{flow+} (Phase~1 with a Phase-2 fallback that
relocates stuck parcels to high-traffic hubs). With 20\% volunteer participation
all Phase-2 heuristics achieve $\geq99\%$ delivery rate for all parcel modes.
\texttt{flow+} is the best overall heuristic.  Delivery rate is approximately
independent of parcel load; the binding constraint is volunteer route coverage,
not throughput.
\end{abstract}

\tableofcontents
\newpage

%----------------------------------------------------------------------
\section{Introduction}
%----------------------------------------------------------------------

Last-mile delivery---the final segment from a distribution hub to the
recipient's door---accounts for a disproportionate share of logistics cost and
urban congestion. Crowdsourced approaches recruit ordinary citizens as
occasional couriers, exploiting the fact that urban residents routinely make
trips that partially overlap with potential parcel routes.

\textbf{Volunteer-based Last-Mile Delivery (VLMD)} is a variant tailored to
metro networks: parcels are placed in station lockers and commuters carry them
on their regular trips. No detour is required---a volunteer simply picks up a
parcel at their boarding station (or a transfer point) if doing so moves the
parcel closer to its destination, and deposits it at the most advanced point
their route reaches. Multiple volunteers relay the parcel across the network
over time.

This paper describes \texttt{vlmd-sim}, a discrete-event simulator for VLMD,
and reports on a systematic experiment using the Porto Metro (Portugal) as a
case study. The main questions are:

\begin{itemize}
  \item What participation rate is needed to achieve a target delivery rate?
  \item How does delivery span evolve with participation and parcel OD mode?
  \item Do flow-weighted heuristics outperform hop-count heuristics?
  \item Is a Phase-2 fallback (hub relocation) beneficial, and under what
        conditions?
\end{itemize}

%----------------------------------------------------------------------
\section{Simulation Model}
%----------------------------------------------------------------------

\subsection{Network}

The network is a directed graph $G=(V,E)$ where $V$ is the set of metro
stations and $E$ contains one directed edge per ordered pair of adjacent
stations on each metro line. The Porto Metro instance has $|V|=85$ stations
and 6 lines. Each edge carries a travel-time attribute (in minutes) and a line
label. A station is a \emph{transfer station} if it is served by more than one
line.

All-pairs unweighted (hop) distances are precomputed at startup via BFS from
every station. Flow-weighted distances (arc weight $= 1/f_{uv}$, where $f_{uv}$
is the total daily OD flow on arc $u\to v$) are computed via Dijkstra.

\subsection{OD Matrix and Volunteers}

The passenger OD matrix is generated from a gravity model applied to
OpenStreetMap land use within 1{,}000~m of each station:
\[
  f_{ij} = \frac{P_i \cdot A_j}{d_{ij}^{1.5}},
\]
where $P_i$ (production) counts residential buildings near station~$i$ and
$A_j$ (attraction) counts commercial, office, and industrial buildings near
station~$j$. Evening flows (hours 17--19) are obtained by reversing morning
flows. The matrix is balanced via iterative proportional fitting (Furness) and
normalised to match Metro do Porto's reported 259{,}000 daily validations
(2025), yielding approximately 195{,}324 OD trips/day.

For each hourly OD entry $(i,j,\mu)$, volunteer counts for hour $h$ of day $d$
are drawn independently from $\mathrm{Poisson}(\mu)$; a fraction $p$ of the
realised commuters become volunteers. Each volunteer travels the shortest
(unweighted) path from $i$ to $j$, with departure times aligned to the GTFS
timetable.

\subsection{Parcels}

Parcel OD pairs are drawn from one of three modes:
\begin{description}
  \item[Uniform:] origin and destination drawn uniformly from $V\times V$
        (excluding self-loops).
  \item[Positive:] OD pair $(i,j)$ sampled with probability proportional to
        $f_{ij}$ (parcels travel where passengers travel---easy case).
  \item[Negative:] OD pair $(i,j)$ sampled with probability proportional to
        $1/f_{ij}$ (parcels travel against the flow---hard case).
\end{description}

Parcel release times are drawn uniformly over the first day of the simulation.
Each instance contains approximately 1{,}950 parcels ($\approx1\%$ of realised
total trips).

\subsection{Locker Interactions}

Volunteers interact with lockers \emph{only at valid stations}: their origin,
their destination, and any line-transfer point along their route. This models
realistic commuter behaviour---no detour is ever required.

\subsection{Transshipment Heuristics}

\subsubsection{Phase 1 (all heuristics)}

Let $d(u,v)$ denote the distance from station $u$ to station $v$ under the
active metric (hop count or flow-weighted). A volunteer at station $s$ carrying
parcel $p$ (destination $t$) drops it off if
\[
  d(s,t) \;\leq\; \min_{s'\in \text{valid ahead}} d(s',t),
\]
i.e.\ no valid station ahead strictly improves on $s$. The parcel is marked
delivered if $d(s,t)=0$.

A volunteer at $s$ not currently carrying a parcel picks up parcel $p$ from
the locker if
\[
  \min_{s'\in \text{valid ahead}} d(s',t_p) \;<\; d(s,t_p),
\]
i.e.\ some valid station ahead is strictly closer to $p$'s destination. When
multiple parcels qualify, the one with the smallest $d(s',t_p)$ (best
reachable distance) is preferred. Phase~1 guarantees monotone progress: a
parcel never moves farther from its destination.

\subsubsection{Phase 2 fallback (\texttt{hop+} / \texttt{flow+})}

Phase~1 can leave parcels stranded at low-traffic stations where no volunteer
route makes progress. Phase~2 activates when Phase~1 finds no candidate and
the following conditions hold:
\begin{enumerate}
  \item The parcel has been waiting in its current locker for at least 1 day
        (1{,}440 minutes).
  \item There exists a valid station $s'$ ahead with higher \emph{combined hub
        attractivity} than the current station:
        \[
          h(s) = \underbrace{F^{\mathrm{OD}}(s)}_{\text{OD flow}} +
                 \underbrace{F^{\mathrm{tr}}(s)}_{\text{transit flow}},
        \]
        where $F^{\mathrm{OD}}(s)$ is the total daily OD flow originating or
        terminating at $s$, and $F^{\mathrm{tr}}(s)$ is the total OD flow that
        passes through $s$ as an intermediate stop on BFS shortest paths.
\end{enumerate}
The 1-day wait threshold is critical: without it, Phase~2 displaces parcels
before same-day Phase~1 volunteers have had a chance to pick them up.

Among eligible parcels the volunteer picks the one closest (under the active
metric) to its destination, maximising the chance of onward delivery.

\subsubsection{Summary of heuristics}

\begin{center}
\begin{tabular}{llcc}
\toprule
Heuristic & Distance metric & Phase 2 & Hub attractivity \\
\midrule
\texttt{hop}   & BFS hop count   & No  & --- \\
\texttt{flow}  & 1/OD-flow       & No  & --- \\
\texttt{hop+}  & BFS hop count   & Yes & $F^{\mathrm{OD}}+F^{\mathrm{tr}}$ \\
\texttt{flow+} & 1/OD-flow       & Yes & $F^{\mathrm{OD}}+F^{\mathrm{tr}}$ \\
\bottomrule
\end{tabular}
\end{center}

\subsection{Discrete-Event Engine}

The simulator maintains a min-heap of events ordered by time. Three event
types are processed:
\begin{description}
  \item[\texttt{parcel\_released}:] place the parcel in its origin locker.
  \item[\texttt{volunteer\_at}:] process locker interaction at the current
        station, then schedule the volunteer's arrival at the next station.
  \item[\texttt{end}:] terminate the simulation (optional hard cutoff).
\end{description}
The simulation runs for 7 days ($10{,}080$ minutes) or until all parcels are
delivered.

%----------------------------------------------------------------------
\section{Experimental Setup}
%----------------------------------------------------------------------

\subsection{Instances}

Ten independent instances are generated by drawing volunteer counts from
$\mathrm{Poisson}(\mu)$ for each OD entry $\mu$. Each instance produces three
parcel files (uniform, positive, negative).

\subsection{Design}

\[
  10 \text{ instances} \;\times\; 3 \text{ parcel modes} \;\times\;
  10 \text{ volunteer fractions} \;\times\; 4 \text{ heuristics}
  \;=\; 1{,}200 \text{ runs}.
\]

Volunteer fractions are log-spaced: 0.1\%, 0.2\%, 0.5\%, 1\%, 2\%, 5\%, 10\%,
20\%, 50\%, 100\%. This range covers realistic crowdsourcing participation rates
(typically 1--5\%) and the saturation region.

A common-random-numbers design is used: for a given (instance, mode) pair the
random seed is fixed, so the volunteer set at fraction $p$ is a strict subset
of those at $p'>p$. Results across fractions and heuristics are therefore
directly comparable.

\subsection{Metrics}

For each run we record:
\begin{itemize}
  \item \textbf{Delivery rate}: fraction of parcels delivered by the end of
        the simulation.
  \item \textbf{Same-day rate}: fraction of parcels delivered on the day of
        release.
  \item \textbf{Delivery span}: time from parcel release to delivery (minutes),
        summarised as mean, median, standard deviation, min, and max.
\end{itemize}

%----------------------------------------------------------------------
\section{Results}
%----------------------------------------------------------------------

\subsection{Delivery Rate vs Volunteer Fraction}

Table~\ref{tab:delivery} shows mean delivery rate $\pm$ standard deviation
across 10 instances at selected volunteer fractions.

\begin{table}[ht]
\centering
\caption{Mean delivery rate (\%) by heuristic, parcel mode, and volunteer
         fraction $p$. Values averaged over 10 instances; standard deviation
         in parentheses.}
\label{tab:delivery}
\small
\begin{tabular}{lrrrrrrrr}
\toprule
 & \multicolumn{2}{c}{Uniform} & \multicolumn{2}{c}{Positive} & \multicolumn{2}{c}{Negative} \\
\cmidrule(lr){2-3}\cmidrule(lr){4-5}\cmidrule(lr){6-7}
$p$ (\%) & \texttt{hop} & \texttt{flow+} & \texttt{hop} & \texttt{flow+} & \texttt{hop} & \texttt{flow+} \\
\midrule
  0.1  & 0.0 (0.0) & 0.0 (0.0) &  0.1 (0.1) &  0.1 (0.1) & 0.0 (0.0) & 0.0 (0.0) \\
  0.5  & 1.6 (0.3) & 1.6 (0.3) & 35.2 (0.9) & 35.5 (1.0) & 0.0 (0.0) & 0.0 (0.1) \\
  1    & 5.3 (0.7) & 5.4 (0.6) & 58.7 (1.3) & 59.2 (1.3) & 0.4 (0.1) & 0.4 (0.1) \\
  2    & 9.5 (0.9) & 9.9 (0.8) & 71.8 (0.9) & 72.0 (0.9) & 1.3 (0.3) & 1.6 (0.3) \\
  5    &26.9 (0.9) &29.1 (0.9) & 88.7 (0.5) & 89.0 (0.5) & 9.7 (1.1) &12.4 (1.1) \\
 10    &58.3 (1.2) &63.7 (1.3) & 98.6 (0.3) & 99.0 (0.3) &38.8 (1.7) &44.5 (1.6) \\
 20    &99.5 (0.4) &100.0(0.1) &100.0 (0.0) &100.0 (0.0) &99.3 (0.5) &100.0(0.1) \\
 50    &99.9 (0.2) &100.0(0.0) &100.0 (0.0) &100.0 (0.0) &99.9 (0.1) &100.0(0.0) \\
100    &100.0(0.0) &100.0(0.0) &100.0 (0.0) &100.0 (0.0) &100.0(0.0) &100.0(0.0) \\
\bottomrule
\end{tabular}
\end{table}

Several patterns are clear:

\begin{itemize}
  \item \textbf{Positive mode} is easiest: 5\% participation achieves
        $\sim$89\% delivery; 10\% achieves $\sim$99\%.
  \item \textbf{Uniform mode} requires $\sim$20\% for near-complete delivery.
  \item \textbf{Negative mode} is hardest: base heuristics need $\sim$20\%
        for 99\% delivery, but exhibit a sharp transition between 10\% and 20\%.
        Phase-2 heuristics (\texttt{flow+}) reach 44.5\% at 10\% vs.\
        38.8\% for \texttt{hop}---a 5.7~pp gain in the most important operating
        range.
  \item At 20\%, \texttt{flow+} achieves 100\% delivery across all modes,
        while \texttt{hop} leaves $\sim$0.7\% undelivered (uniform + negative
        combined).
\end{itemize}

Table~\ref{tab:allheuristics} compares all four heuristics at the sensitive
10\% and 20\% operating points.

\begin{table}[ht]
\centering
\caption{Delivery rate (\%) for all four heuristics at 10\% and 20\%
         volunteer fraction. Values are means over 10 instances.}
\label{tab:allheuristics}
\small
\begin{tabular}{lrrrrrr}
\toprule
 & \multicolumn{3}{c}{$p=10\%$} & \multicolumn{3}{c}{$p=20\%$} \\
\cmidrule(lr){2-4}\cmidrule(lr){5-7}
Heuristic & Uniform & Positive & Negative & Uniform & Positive & Negative \\
\midrule
\texttt{hop}   & 58.3 & 98.6 & 38.8 & 99.5  & 100.0 & 99.3 \\
\texttt{flow}  & 59.5 & 98.0 & 40.9 & 99.5  & 100.0 & 98.7 \\
\texttt{hop+}  & 60.4 & 99.2 & 40.3 & 100.0 & 100.0 & 100.0 \\
\texttt{flow+} & 63.7 & 99.0 & 44.5 & 100.0 & 100.0 & 100.0 \\
\bottomrule
\end{tabular}
\end{table}

\texttt{flow+} leads in uniform and negative modes; \texttt{hop+} leads in
positive mode at 10\%. At 20\%, both Phase-2 heuristics reach 100\% while base
\texttt{hop}/\texttt{flow} fall short by $\sim$0.5--1.3~pp.

\subsection{Same-Day Delivery Rate}

Table~\ref{tab:sameday} reports same-day delivery rates. Since Phase~2 never
fires on day~1 (the 1-day wait threshold ensures this), \texttt{hop+} and
\texttt{hop} have identical same-day rates, as do \texttt{flow+} and
\texttt{flow}. Same-day delivery is substantially harder than overall delivery:
even at 100\% participation, only 97.9\% of uniform parcels are delivered on
the day of release.

\begin{table}[ht]
\centering
\caption{Same-day delivery rate (\%) by heuristic and parcel mode.
         \texttt{hop}/\texttt{hop+} and \texttt{flow}/\texttt{flow+} are
         identical (Phase~2 never fires on day~1).}
\label{tab:sameday}
\small
\begin{tabular}{lrrrrrr}
\toprule
 & \multicolumn{2}{c}{Uniform} & \multicolumn{2}{c}{Positive} & \multicolumn{2}{c}{Negative} \\
\cmidrule(lr){2-3}\cmidrule(lr){4-5}\cmidrule(lr){6-7}
$p$ (\%) & \texttt{hop} & \texttt{flow} & \texttt{hop} & \texttt{flow} & \texttt{hop} & \texttt{flow} \\
\midrule
  5    & 12.7 & 13.2 & 75.6 & 75.8 &  2.5 &  3.0 \\
 10    & 26.4 & 28.2 & 88.3 & 88.5 & 10.8 & 12.5 \\
 20    & 73.5 & 73.2 & 98.2 & 98.0 & 57.6 & 58.7 \\
 50    & 92.7 & 90.7 & 99.6 & 99.6 & 89.2 & 86.2 \\
100    & 97.9 & 97.9 & 99.9 & 99.9 & 97.1 & 96.3 \\
\bottomrule
\end{tabular}
\end{table}

\subsection{Delivery Span}

Table~\ref{tab:span} reports mean delivery span (in minutes) at selected
fractions. Delivery span drops sharply between 10\% and 20\% (the transition
region) and flattens beyond 50\%.

\begin{table}[ht]
\centering
\caption{Mean delivery span (minutes) by heuristic and parcel mode.}
\label{tab:span}
\small
\begin{tabular}{lrrrrrrrr}
\toprule
 & \multicolumn{2}{c}{Uniform} & \multicolumn{2}{c}{Positive} & \multicolumn{2}{c}{Negative} \\
\cmidrule(lr){2-3}\cmidrule(lr){4-5}\cmidrule(lr){6-7}
$p$ (\%) & \texttt{hop} & \texttt{flow+} & \texttt{hop} & \texttt{flow+} & \texttt{hop} & \texttt{flow+} \\
\midrule
  5  & 2024 & 2052 & 471 & 463 & 2963 & 2949 \\
 10  & 2216 & 2363 & 361 & 375 & 3101 & 3166 \\
 20  &  641 &  658 &  71 &  71 &  971 &  998 \\
 50  &  224 &  238 &  35 &  34 &  306 &  329 \\
100  &  114 &  111 &  23 &  24 &  148 &  150 \\
\bottomrule
\end{tabular}
\end{table}

Phase-2 heuristics do not improve span at the same volunteer fraction; they
improve delivery rate (more parcels eventually delivered), which can slightly
inflate reported span by including previously-stuck parcels that required
multiple relay days. At 20\%, where all heuristics deliver nearly 100\%,
\texttt{hop+} achieves the shortest span (597~min uniform, 918~min negative),
while \texttt{flow+} is $\sim$10--8\% longer---the flow-weighted metric routes
relays through high-traffic corridors that may not be the shortest paths for
final delivery.

\subsection{Independence from Parcel Load}

An additional experiment (not detailed here) varied parcel load between
$\sim$195 and $\sim$1{,}950 parcels per instance. Delivery rate was
approximately constant across this range. The binding constraint is volunteer
route topology---can the volunteer relay chain reach the destination?---not
throughput capacity. The 7-day window drains any queue that forms at bottleneck
stations.

\subsection{Failure Mode Analysis}

Undelivered parcels fall into two categories:
\begin{description}
  \item[At-origin (never picked up):] the dominant failure mode. Root causes:
    \begin{itemize}
      \item No volunteer ever visits the origin station.
      \item All volunteers whose route passes through the origin terminate there
            (no valid station ahead to improve distance).
      \item Volunteers pass through but no onward valid station reduces distance
            to destination.
    \end{itemize}
  \item[In-transit (stranded mid-route):] relay chain broke partway. Common at
        terminus stations due to directional flow imbalance.
\end{description}

\paragraph{Terminus traps.}
Stations at the ends of lines (e.g.\ Vila do Conde, Varziela, Ramalde in the
Porto Metro) accumulate undeliverable parcels because the dominant volunteer
flow is city-inbound during the morning. Outbound parcels are structurally
trapped: Phase~2 can move parcels to intermediate high-traffic hubs (e.g.\
Senhora da Hora) but cannot complete the final outbound leg because no
outbound-bound volunteer carries them further. This is a network topology
limitation not addressable by heuristic changes alone.

Under the negative parcel mode at 10\% volunteers, peak locker occupancy at
terminus stations reaches 138 parcels (Vila do Conde), with a mean of $\sim$106
parcels over the 7-day window---confirming that parcels arrive but rarely leave.

\subsection{Design of the Phase-2 Fallback}

A series of heuristic variants was tested to disentangle the key design choices:

\begin{description}
  \item[\texttt{hop++}/\texttt{flow++}:] time-of-day threshold (fallback fires
        noon--18:00 to transfer stations, after 18:00 to any busier station),
        using transit flow as hub attractivity. Tested and retired: achieves
        $\sim$1~pp higher delivery rate at 10\% negative but loses $\sim$7~pp
        same-day rate at 100\% volunteers.
  \item[\texttt{hop+++}/\texttt{flow+++}:] crossing variant combining transit
        flow with the 1-day threshold. Tested and retired: indistinguishable
        from \texttt{hop+}/\texttt{flow+} in all metrics, confirming that the
        \emph{threshold} (1-day wait) is the dominant design parameter, not
        the hub-selection metric.
\end{description}

The final design uses \emph{combined hub attractivity}
$h(s)=F^{\mathrm{OD}}(s)+F^{\mathrm{tr}}(s)$ for robustness, even though
empirically $F^{\mathrm{OD}}$ alone would yield the same results for Porto
Metro (major interchange stations dominate both signals).

%----------------------------------------------------------------------
\section{Conclusion}
%----------------------------------------------------------------------

We presented \texttt{vlmd-sim}, a discrete-event simulator for volunteer-based
last-mile delivery over metro networks, and reported results from a 1{,}200-run
experiment on the Porto Metro.

\textbf{Main findings:}

\begin{enumerate}
  \item \textbf{20\% participation is sufficient} to achieve $\geq99\%$
        delivery rate for all parcel modes, using Phase-2 heuristics.
        The transition from low to high delivery rate is sharp and occurs
        between 10\% and 20\%.

  \item \textbf{\texttt{flow+} is the best overall heuristic.} The flow-weighted
        distance metric helps most in the negative parcel mode (5--20\%
        fractions), where parcels travel against the dominant passenger flow.
        Phase~2 provides an additional 3--6~pp at 10\% negative.

  \item \textbf{Same-day delivery requires much higher participation.}
        At 10\% volunteers, only $\sim$26--28\% of uniform parcels are delivered
        on the day of release. Phase~2 does not reduce same-day rate because its
        1-day threshold leaves day-1 Phase-1 opportunities intact.

  \item \textbf{Delivery rate is independent of parcel load.}
        The system bottleneck is volunteer route coverage, not throughput.

  \item \textbf{Terminus traps are the primary structural limitation.}
        Stations at the ends of lines accumulate undeliverable parcels.
        Addressing this requires either demand-side incentives (rewarding
        outbound carries) or supply-side infrastructure (shuttle or last-mile
        service from terminus hubs).

  \item \textbf{The 1-day wait threshold is the critical Phase-2 parameter.}
        The hub-selection metric (OD flow vs.\ transit flow, single-tier vs.\
        two-tier) has negligible impact on outcomes. The threshold dominates.
\end{enumerate}

\textbf{Future work} includes: (i) explicit locker capacity constraints;
(ii) per-volunteer acceptance probability $q < 1$; (iii) non-shortest-path
volunteer routing; (iv) timetable-driven inter-station travel times replacing
the fixed 3-minute leg fallback; and (v) unit tests for the network and
scheduler modules.

\bigskip
\noindent The simulator source code is available at:
\begin{center}
\url{https://github.com/joaopedroso/vlmd-sim}
\end{center}

\end{document}
